\documentclass[9.5pt,a4paper]{article}
\usepackage[margin=13mm,top=10mm,bottom=10mm]{geometry}
\usepackage{fontspec}
\usepackage[french]{babel}
\usepackage{xcolor}
\usepackage[most]{tcolorbox}
\usepackage{booktabs}
\usepackage{enumitem}
\usepackage{graphicx}
\usepackage{array}
\usepackage{titlesec}
\usepackage{listings}

\definecolor{bleu}{HTML}{1F6FEB}
\definecolor{gris}{HTML}{5B6470}
\definecolor{vert}{HTML}{1E8449}
\definecolor{orange}{HTML}{B45309}
\definecolor{fond}{HTML}{F4F6F8}
\definecolor{codebg}{HTML}{F6F8FA}

\setlength{\parindent}{0pt}
\titlespacing*{\section}{0pt}{7pt}{2pt}
\titleformat{\section}{\normalfont\bfseries\color{bleu}\normalsize}{\thesection}{0.6em}{}
\setlist[itemize]{nosep,leftmargin=1.05em,labelsep=.4em}
\setlist[enumerate]{nosep,leftmargin=1.25em,labelsep=.45em}

\lstdefinestyle{py}{language=Python,backgroundcolor=\color{codebg},
  basicstyle=\ttfamily\fontsize{7.4}{8.6}\selectfont,
  keywordstyle=\color{bleu}\bfseries,commentstyle=\color{vert}\itshape,
  stringstyle=\color{orange},frame=none,breaklines=true,
  showstringspaces=false,columns=fullflexible,keepspaces=true,
  aboveskip=2pt,belowskip=2pt,
  literate={é}{{\'e}}1 {è}{{\`e}}1 {à}{{\`a}}1 {ç}{{\c c}}1 {û}{{\^u}}1
           {→}{{$\rightarrow$}}1 {≥}{{$\geq$}}1}
\lstset{style=py}

\newtcolorbox{piege}[1]{enhanced,colback=orange!6,colframe=orange!80!black,boxrule=.5pt,
  arc=2pt,left=4pt,right=4pt,top=2pt,bottom=2pt,
  fonttitle=\bfseries\footnotesize,title={#1}}
\newtcolorbox{astuce}[1]{enhanced,colback=vert!5,colframe=vert!70!black,boxrule=.5pt,
  arc=2pt,left=4pt,right=4pt,top=2pt,bottom=2pt,
  fonttitle=\bfseries\footnotesize,title={#1}}

\begin{document}

\begin{tcolorbox}[enhanced,colback=fond,colframe=fond,arc=3pt,
  left=8pt,right=8pt,top=5pt,bottom=5pt]
\begin{minipage}[c]{0.78\textwidth}
  {\color{gris}\footnotesize LICENCE 3 SDIA $\cdot$ ENSPD --- UNIVERSITÉ DE DOUALA}\\[2pt]
  {\Large\bfseries Python pour le machine learning --- fiche de révision}\\[3pt]
  {\small L'essentiel pour les TP : bases, NumPy, pandas, matplotlib, scikit-learn}\\[3pt]
  {\color{gris}\footnotesize Dr.\ Noulapeu Ngaffo Armielle}
\end{minipage}\hfill
\begin{minipage}[c]{0.18\textwidth}
  \raggedleft\includegraphics[width=16mm]{logo_enspd.png}
\end{minipage}
\end{tcolorbox}

\section*{\color{bleu}1. Les bases, en dix lignes}
\begin{lstlisting}
x = 12                 # entier
prix = 12.5            # reel
nom = "Ada"            # chaine de caracteres
ok = True              # booleen
notes = [12, 15, 9]    # liste
etudiant = {"nom": "Ada", "note": 15}   # dictionnaire
notes[0]               # 12
etudiant["note"]       # 15
for n in notes:        # boucle
    print(n)
if x > 10:             # condition
    print("grand")
def moyenne(liste):    # fonction
    return sum(liste) / len(liste)
\end{lstlisting}

\begin{minipage}[t]{0.49\textwidth}
\begin{astuce}{Ce que vous écrirez vraiment}
\begin{itemize}
  \item une liste pour les colonnes à garder ;
  \item un dictionnaire pour les réglages ;
  \item une fonction pour éviter de copier-coller ;
  \item une boucle pour parcourir des fichiers.
\end{itemize}
\end{astuce}
\end{minipage}\hfill
\begin{minipage}[t]{0.49\textwidth}
\begin{piege}{Les trois erreurs de débutant}
\textbf{IndentationError} : Python lit l'indentation (4 espaces, jamais de
tabulation).\\
\textbf{IndexError} : \texttt{notes[3]} n'existe pas dans une liste de 3.\\
\textbf{KeyError} : la clé du dictionnaire est mal orthographiée.
\end{piege}
\end{minipage}

\section*{\color{bleu}2. NumPy : le tableau, brique de base}
\begin{lstlisting}
import numpy as np
a = np.array([[1, 2, 3], [4, 5, 6]])
a.shape          # (2, 3)  -> 2 lignes, 3 colonnes
a.mean(axis=0)   # [2.5 3.5 4.5]   moyenne par colonne
a.mean(axis=1)   # [2. 5.]          moyenne par ligne
a[:, 1]          # [2 5]            deuxieme colonne
a[a > 3]         # [4 5 6]          filtre booleen
a + 10           # ajoute 10 a chaque case
\end{lstlisting}

\begin{astuce}{La règle de l'axe, à retenir une fois pour toutes}
\texttt{axis=0} écrase les \textbf{lignes} : on obtient une valeur par
\textbf{colonne}. \quad \texttt{axis=1} écrase les \textbf{colonnes} : une
valeur par \textbf{ligne}. C'est la source d'erreur numéro un.
\end{astuce}

\begin{piege}{Formes}
\texttt{a.shape} doit correspondre à ce qu'attend la fonction suivante.
Une erreur \texttt{shapes not aligned} vient presque toujours d'un désaccord
de dimensions, pas d'un bug de la bibliothèque.
\end{piege}

\section*{\color{bleu}3. pandas : le tableau qui a des noms}
\begin{lstlisting}
import pandas as pd
df = pd.read_csv("donnees.csv")
df.head()                  # 5 premieres lignes
df.shape                   # (lignes, colonnes)
df.info()                  # types et valeurs manquantes
df.describe()              # min, max, moyenne, quartiles
df["note"]                 # une colonne (Series)
df[["note", "absences"]]   # plusieurs colonnes
df[df["note"] > 12]        # filtre
df.loc[0, "note"]          # ligne 0, colonne "note"
df.iloc[0, 1]              # ligne 0, colonne 1 (position)
df.groupby("groupe")["note"].mean()
df.isna().sum()            # manquants par colonne
df["note"].fillna(df["note"].median())
\end{lstlisting}

\begin{astuce}{Sur un tableau de 4 lignes (vérifié)}
\texttt{groupe = A, A, B, B} ; \texttt{note = 12, 15, 9, 13}.\\
\texttt{df[df.note > 12]} garde les notes \textbf{15 et 13}.\\
\texttt{groupby("groupe").mean()} donne \textbf{A = 13,5} et \textbf{B = 11,0}.\\
\texttt{isna().sum()} vaut \textbf{0} : aucune valeur manquante.
\end{astuce}

\begin{piege}{Deux pièges classiques}
\textbf{SettingWithCopy} : ne modifiez pas une tranche ; utilisez
\texttt{df.loc[condition, "colonne"] = valeur}.\\
\textbf{Index désaligné} : après un filtre, l'index garde des trous ;
\texttt{reset\_index(drop=True)} avant de concaténer.
\end{piege}

\section*{\color{bleu}4. matplotlib : montrer, pas décorer}
\begin{lstlisting}
import matplotlib.pyplot as plt
plt.figure(figsize=(7, 4))
plt.plot(x, y, label="observe")
plt.scatter(x, y, s=10)
plt.hist(valeurs, bins=20)
plt.xlabel("temps"); plt.ylabel("valeur")
plt.title("Titre qui dit ce qu'on voit")
plt.legend(); plt.show()
\end{lstlisting}

\begin{astuce}{Les quatre obligations d'une figure}
\begin{enumerate}
  \item un titre qui énonce le constat, pas le nom de la variable ;
  \item des axes nommés, avec l'unité ;
  \item une légende dès qu'il y a deux courbes ;
  \item une échelle honnête --- pas d'axe tronqué pour exagérer un écart.
\end{enumerate}
\end{astuce}

\section*{\color{bleu}5. scikit-learn : cinq lignes, toujours les mêmes}
\begin{lstlisting}
from sklearn.model_selection import train_test_split
from sklearn.pipeline import Pipeline
from sklearn.preprocessing import StandardScaler
from sklearn.linear_model import LogisticRegression
from sklearn.metrics import accuracy_score

X_train, X_test, y_train, y_test = train_test_split(
    X, y, test_size=0.25, random_state=4112, stratify=y)

modele = Pipeline([("sc", StandardScaler()),
                   ("clf", LogisticRegression(max_iter=1000))])
modele.fit(X_train, y_train)      # apprendre
pred = modele.predict(X_test)     # predire
accuracy_score(y_test, pred)      # mesurer
\end{lstlisting}

\begin{astuce}{Le vocabulaire de l'API}
\texttt{fit} : apprendre sur l'entraînement. \quad
\texttt{predict} : produire une réponse. \quad
\texttt{score} : mesurer sur des données nouvelles.

Un objet scikit-learn a toujours ces méthodes : changer de modèle ne change
que le nom de la classe.
\end{astuce}

\begin{piege}{Les formes attendues}
\texttt{X} : matrice $(n_{\text{exemples}}, n_{\text{variables}})$ --- ici
\texttt{(4, 2)}. \quad \texttt{y} : vecteur de longueur $n$.
Passer une colonne au lieu d'une matrice provoque
\texttt{Expected 2D array}.
\end{piege}

\section*{\color{bleu}6. Diagnostiquer une erreur, dans l'ordre}
\begin{center}
\footnotesize
\begin{tabular}{@{}p{0.27\textwidth}p{0.30\textwidth}p{0.33\textwidth}@{}}
\toprule
\textbf{Message} & \textbf{Cause probable} & \textbf{Correction} \\
\midrule
\texttt{ModuleNotFoundError} & paquet absent ou mauvais noyau &
\texttt{pip install}, vérifier le noyau Jupyter \\
\texttt{FileNotFoundError} & chemin relatif erroné &
afficher \texttt{os.getcwd()}, utiliser un chemin absolu \\
\texttt{KeyError} & nom de colonne mal écrit & afficher \texttt{df.columns} \\
\texttt{ValueError: NaN} & valeurs manquantes & imputer ou retirer \\
\texttt{Expected 2D array} & \texttt{X} est une Series & écrire
\texttt{X[["col"]]} \\
\texttt{shapes not aligned} & dimensions incompatibles & afficher les
\texttt{.shape} \\
\bottomrule
\end{tabular}
\end{center}

\begin{astuce}{Réflexe de débogage}
Affichez \textbf{avant} de calculer : \texttt{type(x)}, \texttt{x.shape},
\texttt{x.head()}. Quatre-vingts pour cent des erreurs se voient là, avant
même d'être levées.
\end{astuce}

\section*{\color{bleu}7. Six exercices de vérification}
\begin{enumerate}
  \item Créer une liste de 5 notes et en calculer la moyenne sans NumPy.
  \item Créer un tableau $3 \times 4$ et calculer la somme de chaque colonne.
  \item Charger un CSV et afficher le nombre de valeurs manquantes par colonne.
  \item Filtrer les lignes où une note dépasse 10, puis compter les lignes.
  \item Calculer la moyenne d'une colonne par groupe.
  \item Découper un jeu en entraînement et test, entraîner une régression
        logistique et afficher l'exactitude.
\end{enumerate}

\vspace{1mm}
{\footnotesize\color{gris}Si l'un de ces six gestes bloque, c'est celui-là
qu'il faut revoir avant le prochain TP : ils couvrent l'essentiel de ce que
vous aurez à écrire.}

\end{document}
